\chapter{Model--Specification Correspondence}
\label{ch:model-spec-correspondence}

A machine-checked security proof is only as persuasive as the correspondence
between its formal model and the scheme that the proof purports to study.  The
HAETAE proof surface therefore has two distinct obligations.  The first is the
internal EasyCrypt obligation: the games, hops, and reductions must check.  The
second is the external modeling obligation: the objects in those games must
faithfully represent the HAETAE specification.  This chapter records that second
obligation explicitly.

The specification reference used here is \ec{HAETAE\_v260204.pdf}
\cite{haetae-v260204}.  The local copy has 46 pages and was generated on
February 4, 2026.  The relevant parts for provable security are Definition 1 for
the signature-scheme interface, Section 4.2 for the HAETAE algorithms, Figure 8
for the deterministic full algorithms, and Table 1 for the three parameter
sets.  The EasyCrypt files used for the
formal model are principally \ec{Sig\_ROM.ec}, \ec{HAETAE\_Scheme.ec},
\ec{HAETAE\_Params.ec}, \ec{HAETAE\_ROM.ec}, \ec{HAETAE\_Distributions.ec},
\ec{HAETAE\_Algebra.ec}, and \ec{HAETAE\_Rejection.ec}.

\section{Signature-scheme interface}
\label{sec:model-spec-signature-interface}

Definition 1 of the specification presents a signature scheme as a triple of
algorithms
\[
  \Pi=(\mathsf{KeyGen},\mathsf{Sign},\mathsf{Verify})
\]
with public keys, secret keys, messages, signatures, and a verification
predicate.  The corresponding EasyCrypt interface is intentionally generic.  It
first defines abstract types and then requires any concrete scheme to implement
key generation, signing, and verification procedures.

\begin{longtable}{p{0.25\textwidth}p{0.30\textwidth}p{0.35\textwidth}}
\caption{Signature-interface correspondence.}
\label{tab:signature-interface-correspondence}\\
\toprule
\textbf{Specification object} & \textbf{EasyCrypt object} & \textbf{Mathematical reading} \\
\midrule
\endfirsthead
\toprule
\textbf{Specification object} & \textbf{EasyCrypt object} & \textbf{Mathematical reading} \\
\midrule
\endhead
Public key space & \ec{type pkey} in \ec{Sig\_ROM.ec} & A carrier set
\(\mathcal{PK}\). \\
\midrule
Secret key space & \ec{type skey} & A carrier set \(\mathcal{SK}\). \\
\midrule
Message and context & \ec{type message}, \ec{type context} & The security game works over
pairs \((m,ctx)\in\mathcal{M}\times\mathcal{C}\). \\
\midrule
Signature space & \ec{type signature} & A carrier set \(\Sigma\). \\
\midrule
Key generation & \ec{proc kg() : pkey * skey} & A probabilistic kernel
\(\mathsf{KeyGen}:1\to\mathcal{D}(\mathcal{PK}\times\mathcal{SK})\). \\
\midrule
Signing & \ec{proc sign(sk,m,ctx) : signature} & A random-oracle relative
kernel \(\mathsf{Sign}^{H}(sk,m,ctx)\). \\
\midrule
Verification & \ec{proc verify(pk,m,ctx,sig) : bool} & A predicate
\(\mathsf{Verify}^{H}(pk,m,ctx,\sigma)\in\{0,1\}\). \\
\midrule
Chosen-message and strong freshness & \ec{fresh\_msg}, \ec{fresh\_sig} &
\ec{fresh\_msg} is the EUF-CMA condition that the final message/context pair was
not previously signed.  \ec{fresh\_sig} is the stronger condition that the exact
message/context/signature triple was not returned before. \\
\bottomrule
\end{longtable}

The generic interface is not itself a HAETAE model.  It is the mathematical
ambient category in which the HAETAE model is later instantiated.  In categorical
language, \ec{HAETAE\_Scheme.ec} supplies one object of the class of all
random-oracle signature schemes described by \ec{Sig\_ROM.ec}.

\section{Algorithm correspondence}
\label{sec:model-spec-algorithm-correspondence}

Figure 8 of the specification gives deterministic HAETAE algorithms.  The
EasyCrypt wrapper in \ec{HAETAE\_Scheme.ec} exposes the same three algorithms,
but separates the random-oracle queries from the internal algebraic operations.
This separation is essential for the game proof: the ROM table, query logs, and
programming events must be visible to EasyCrypt.

\begin{longtable}{p{0.23\textwidth}p{0.34\textwidth}p{0.33\textwidth}}
\caption{Algorithm-level model/spec correspondence.}
\label{tab:algorithm-correspondence}\\
\toprule
\textbf{Specification step} & \textbf{EasyCrypt representation} & \textbf{Proof-relevant interpretation} \\
\midrule
\endfirsthead
\toprule
\textbf{Specification step} & \textbf{EasyCrypt representation} & \textbf{Proof-relevant interpretation} \\
\midrule
\endhead
Sample key-generation seed and derive \(seed_A,seed_{sk},K\). &
\ec{sd <\$ dseed}; \ec{haetae\_keygen\_rhoprime sd}; \ec{keygen\_internal}. &
The model records the entropy source and delegates deterministic key material
construction to the internal HAETAE operations. \\
\midrule
Expand the public matrix from \(seed_A\). &
\ec{H.get(matrix\_expand\_query haetae\_mode rhoprime)} and
\ec{keygen\_internal}. &
The wrapper records a typed matrix-expansion ROM query, while
\ec{keygen\_internal} constructs the modeled matrix through deterministic
internal algebra.  The ROM output is not consumed by \ec{keygen\_internal}; it is
part of the proof instrumentation for the hash-interface audit. \\
\midrule
Compute \(b=A_{gen}s_{gen}+e_{gen}+a_{gen}\), split low/high bits, and form
\((pk,sk)\). &
\ec{keygen\_internal haetae\_mode rhoprime}. &
The algebraic construction is part of the scheme model, not part of the ROM
wrapper.  Its soundness depends on the algebra and parameter files. \\
\midrule
Compute \(\mu=H_{gen}(seed_A,b_1,M)\). &
\ec{H.get(message\_hash\_query pk ctx m)} followed by \ec{ro\_message\_hash}. &
The signing wrapper performs and logs a typed message-hash ROM query.  In the
current internal algebra, \ec{sign\_internal} recomputes the message hash through
the deterministic operation \ec{message\_hash}; the wrapper's ROM output is
therefore observational instrumentation unless the model is changed to pass it
into \ec{sign\_internal}. \\
\midrule
Derive signing randomness through \(seed_{ybb}=H_{gen}(K,\mu)\) and
\(expandYbb\). &
\ec{coins <\$ drandom\_coins}; \ec{H.get(sampler\_expand\_query coins)};
\ec{ro\_signing\_coins}. &
The formal proof models the signing sampler as a ROM-mediated randomness
source.  The model/spec audit must justify that this is the intended abstraction
of the deterministic specification path. \\
\midrule
Hash the commitment data to obtain challenge seed \(\rho\), then sample the
binary challenge. &
\ec{commitment\_highbits}, \ec{commitment\_lowbits},
\ec{H.get(challenge\_hash\_query haetae\_mode highbits lowbits mu)}, and
\ec{sign\_internal}. &
The challenge point is a separately typed ROM site for freshness and
programming analysis.  The explicit wrapper call uses the wrapper value
\ec{mu <- ro\_message\_hash ro\_y}.  The signature returned by
\ec{sign\_internal}, and the honest programming site used by the programming
lemmas, instead use the deterministic \ec{message\_hash} recomputed inside the
internal algebra.  Thus this wrapper call records a typed challenge-hash query
near the proof site; it should not be identified with the transcript-derived
programming site unless a separate bridge equates the two message-hash values. \\
\midrule
Apply rejection tests and return encoded signature components. &
\ec{sign\_internal}, \ec{HAETAE\_Rejection.ec}, and distribution lemmas. &
The proof reasons about support, rejection predicates, and exact-hyperball
sampling through explicit events and loss lemmas. \\
\midrule
Decode, recompute challenge data, and check norm/challenge conditions. &
\ec{verify\_internal haetae\_mode pk m ctx sig}. &
Verification is modeled as the deterministic predicate used by the correctness
and EUF-CMA games. \\
\bottomrule
\end{longtable}

The most important modeling distinction is the treatment of signing randomness.
The specification describes a deterministic derivation from \(K\) and \(\mu\),
whereas the EasyCrypt signing wrapper first samples \ec{random\_coins} and then
routes the sampler-expansion point through the random oracle.  This is not a
minor notation change.  It is the reason that the proof contains
\ec{sampler\_expand\_query}, \ec{sampler\_bad\_prequery},
\ec{sampler\_rom\_covered}, and the concrete lazy-ROM freshness lemmas.  The
formal theorem is therefore a theorem about the modeled ROM interface; a
complete certification story must additionally maintain a model/spec argument
that the deterministic HAETAE sampling path realizes the same distributional
role.

\section{Exact scope statement}
\label{sec:model-spec-exact-scope}

The model/specification correspondence in this chapter should be read as an
audited mathematical bridge, not as an additional EasyCrypt theorem silently
imported into the proof.  Let
\[
  \mathsf{Spec}_{260204}
\]
denote the algorithms in the February 2026 specification, and let
\[
  \mathsf{ECModel}
\]
denote the EasyCrypt scheme obtained from \ec{HAETAE\_Scheme.ec} and its
supporting files.  The machine-checked EasyCrypt surface proves the following
conditional formal-model implication, where \(\mathcal{H}_{\mathsf{EC}}\)
abbreviates the final hop/reduction premise and the MLWE and Module-SIS
hardness premises appearing in the cited EasyCrypt composition lemma:
\[
  \mathcal{H}_{\mathsf{EC}}
  \Longrightarrow
  \Adv_{\mathsf{ECModel},\rom}^{\eufcma}(A)
  \le
  \ec{haetae\_euf\_counted\_rom\_bound}.
\]
The dissertation does not claim a separately checked theorem of the form
\[
  \mathsf{Spec}_{260204}\equiv\mathsf{ECModel}.
\]
Instead, it records the correspondence evidence needed for a reviewer to judge
whether the formal model is the intended model of the specification.

This distinction matters most for signing randomness.  In the specification,
the signing path derives \(\mathsf{seed}_{ybb}\) from secret key material and the
message digest, then expands it.  In the EasyCrypt proof, the sampler expansion
site is exposed as a typed random-oracle query fed by sampled
\ec{random\_coins}.  The security theorem is exact for the EasyCrypt model; the
specification-level conclusion additionally depends on the distributional
justification that the exposed sampler interface faithfully represents the
specified signing-randomness role.

\begin{assumption}[Signing-randomness modeling bridge]
\label{assm:signing-randomness-modeling-bridge}
The specification-level reading of the EasyCrypt theorem assumes the following
modeling hypotheses.
\begin{enumerate}[label=\arabic*.]
\item The specification's call
\[
  \mathsf{seed}_{ybb}=H_{gen}(K,\mu),
  \qquad
  (y,b,b')=\mathsf{expandYbb}(\mathsf{seed}_{ybb},counter)
\]
is represented in the EasyCrypt model by a domain-separated sampler-expansion
random-oracle site
\[
  q_s=\ec{sampler\_expand\_query}(\rho_s)
\]
and interpretation map \(\ec{ro\_signing\_coins}\).
\item The sampler, message-hash, challenge-hash, and matrix-expansion query
constructors are injectively tagged, so that programming or observing one domain
does not program or observe another.
\item Before the signing call under consideration, \(q_s\) has not appeared in
the adversary hash log or in the sampler self-log from earlier signing calls
except on the local wrapper event \(\ec{sampler\_bad\_prequery}\).  The
counted-ROM prequery and min-entropy terms are separate final-bound ledger
entries, not this local event.
\item At a fresh sampler site, the concrete lazy-ROM law
\[
  \ec{ro\_signing\_coins}(H(q_s)) \equiv D_{\mathsf{coins}}
\]
is the law expressed by
\ec{concrete\_lazy\_rom\_sampler\_expand\_query\_fresh\_eq}.
\end{enumerate}
Under these hypotheses, the distribution of the specification signing sampler,
viewed through the random-oracle abstraction, is identified with the EasyCrypt
sampler distribution on the clean event.  Without conditioning on the clean
event, two distinct proof-accounting facts are involved: the local
\ec{sampler\_bad\_prequery} wrapper event is bounded by the combined
adversary-hash and sampler-self-log budget, while the final counted-ROM theorem
contains the separate \ec{counted\_rom\_prequery\_term} and
\ec{counted\_rom\_min\_entropy\_term} summands.
\end{assumption}

\paragraph{Pen-and-paper justification.}
By domain separation, the sampler site is disjoint from the message, challenge,
and matrix sites.  On the clean event, the sampler site is fresh in the lazy-ROM
table.  The lazy-ROM semantics therefore samples a new oracle output from the
declared output distribution, and the interpretation
\ec{ro\_signing\_coins} maps that output to the same coin distribution used by
the EasyCrypt signing sampler.  The local sampler equivalence can fail before
the fresh call only if the adversary has already queried the sampler site or an
earlier signing sampler has already logged the same sampler site.  These
failures are represented by \(\ec{sampler\_bad\_prequery}\) and have their own
combined-log budget in the active-signing wrapper.  The counted-ROM prequery and
min-entropy terms are separate final-bound summands for the ROM programming
ledger: the former is a programmed-site prequery term, and the latter records
the \(\ec{bad\_min\_entropy}\) predicate on observed transcripts.  Thus the
bridge is exact under \(\neg\ec{sampler\_bad\_prequery}\), with the remaining
accounting split between the local sampler-prequery bound and the counted-ROM
terms.

This assumption is intentionally not presented as an EasyCrypt theorem.  It is
the strongest pen-and-paper bridge currently supplied by the dissertation.
Mechanizing it would require a formal relation between the deterministic
\(H_{gen}(K,\mu)\)-based specification sampler and the typed
\ec{sampler\_expand\_query} interface used by the current proof.

\section{Internal-algorithm fidelity audit ledger}
\label{sec:internal-algorithm-fidelity-ledger}

The previous assumption isolates the signing-randomness abstraction.  The same
discipline is needed for the remaining internal algorithms: a theorem about
\(\mathsf{ECModel}\) is only a theorem about \(\mathsf{Spec}_{260204}\) after
the following correspondences have either been mechanized or accepted as
audited modeling assumptions.

The status column uses the dependency classes from
\Cref{tab:dependency-classes}.  In particular, ``checked for
\(\mathsf{ECModel}\)'' means that the EasyCrypt model has the stated internal
property; it does not mean that the corresponding line of the PDF specification
has been formally related to the EasyCrypt definition.  ``Audited'' means
reviewed correspondence evidence.  ``Assumed'' means a named bridge that remains
outside the checked proof.

\begin{longtable}{p{0.21\textwidth}p{0.23\textwidth}p{0.32\textwidth}p{0.14\textwidth}}
\caption{Internal-algorithm correspondence obligations for reading
\(\mathsf{ECModel}\) as the HAETAE specification.}
\label{tab:internal-algorithm-fidelity-ledger}\\
\toprule
\textbf{Specification object} & \textbf{EasyCrypt object} &
\textbf{Current evidence} & \textbf{Status} \\
\midrule
\endfirsthead
\toprule
\textbf{Specification object} & \textbf{EasyCrypt object} &
\textbf{Current evidence} & \textbf{Status} \\
\midrule
\endhead
Figure 8 KeyGen, lines 1--11 &
\ec{keygen\_internal}, \ec{HAETAE\_Scheme.kg} &
The parameter, matrix-expansion, secret-sampling, and public-key equations are
represented by the algebraic model and the parameter table in
\Cref{sec:model-spec-parameter-correspondence}. &
Audited; not a checked PDF-to-model equivalence. \\
\midrule
Figure 8 Sign, lines 1--17 &
\ec{sign\_internal}, \ec{HAETAE\_Scheme.sign} &
The proof tracks rejection, sampler expansion, message hashing, challenge
hashing, and signature validity through the hop-game lemmas and the clean
sampler instrumentation. &
Checked for \(\mathsf{ECModel}\); specification bridge assumed. \\
\midrule
Figure 8 Verify, lines 1--10 &
\ec{verify\_internal}, \ec{HAETAE\_Scheme.verify} &
The verification predicate is the formal accept/reject event used by the
EUF-CMA game, but byte-level parsing and malformed encodings are outside the
central probability theorem. &
Audited; implementation parsing separate. \\
\midrule
Table 1 parameters &
\ec{Mode2}, \ec{Mode3}, \ec{Mode5}, \ec{mode\_*} &
The dissertation records the main numerical correspondences, including
\(n=256\), \(q=64513\), \((k,\ell)\), \(\tau\), and squared norm bounds. &
Documented correspondence. \\
\midrule
\(H_{gen}(K,\mu)\) and \(\mathsf{expandYbb}\) &
\ec{sampler\_expand\_query}, \ec{ro\_signing\_coins} &
The EasyCrypt proof has concrete lazy-ROM freshness and self-query logging for
the modeled sampler site.  The deterministic specification-to-ROM distribution
argument is exactly \Cref{assm:signing-randomness-modeling-bridge}. &
Explicit assumption. \\
\bottomrule
\end{longtable}

\section{Parameter correspondence}
\label{sec:model-spec-parameter-correspondence}

Table 1 of the specification gives the parameter sets for security levels 2, 3,
and 5.  The EasyCrypt file \ec{HAETAE\_Params.ec} encodes the same modes through
\ec{Mode2}, \ec{Mode3}, and \ec{Mode5}.  The following table records the main
correspondences used by the proof.

\begin{longtable}{p{0.22\textwidth}p{0.28\textwidth}p{0.40\textwidth}}
\caption{Parameter correspondence between the specification and EasyCrypt.}
\label{tab:parameter-correspondence}\\
\toprule
\textbf{Specification parameter} & \textbf{EasyCrypt operation} & \textbf{Values / interpretation} \\
\midrule
\endfirsthead
\toprule
\textbf{Specification parameter} & \textbf{EasyCrypt operation} & \textbf{Values / interpretation} \\
\midrule
\endhead
Security level & \ec{mode} with constructors \ec{Mode2}, \ec{Mode3}, \ec{Mode5} &
The three formal modes correspond to the three specification parameter sets. \\
\midrule
Ring dimension \(n\) & \ec{n} & \(256\). \\
\midrule
Modulus \(q\) & \ec{q} & \(64513\); \ec{dq} is \(2q\). \\
\midrule
Module dimensions \((k,\ell)\) & \ec{mode\_k}, \ec{mode\_l} &
\((2,4)\), \((3,6)\), and \((4,7)\). \\
\midrule
Challenge weight \(\tau\) & \ec{mode\_tau} &
\(58\), \(80\), and \(128\). \\
\midrule
Norm/radius bounds \(B,B',B''\) & \ec{mode\_b0sq}, \ec{mode\_b1sq},
\ec{mode\_b2sq} & The EasyCrypt model stores squared integer bounds.  The
correspondence is therefore between the specification radii and the squared
thresholds used by norm predicates. \\
\midrule
Hint parameters & \ec{mode\_alpha\_hint}, \ec{mode\_m} &
Mode-dependent constants used by encoding, hint, and verification predicates. \\
\midrule
Public-key and signature sizes & \ec{mode\_publickeybytes},
\ec{mode\_crypto\_bytes} & Size constants used for specification alignment and
encoding-side modeling. \\
\bottomrule
\end{longtable}

The squared-bound convention is a common source of confusion.  The main
EasyCrypt validity predicates compare squared norms against squared integer
thresholds with non-strict bounds.  The proof should not silently identify
ordinary norm notation with the exact predicate used in the model; the intended
reading for these validity checks is
\[
  \|z\|^2 \le B^2
  \quad\text{implemented as}\quad
  \mathsf{normsq}(z) \le \ec{mode\_b*sq}(\mathsf{mode}).
\]
Some rejection subchecks use strict inequalities against separate scaled
rejection bounds; those are distinct abort predicates, not the main
\ec{mode\_b*sq} validity checks.

\section{Hash and random-oracle correspondence}
\label{sec:model-spec-hash-correspondence}

The specification uses hash/XOF building blocks \(H_{gen}\), \(H\),
\(expandA\), \(expandS\), and \(expandYbb\).  The EasyCrypt proof distinguishes
three layers:
\[
  \text{typed query constructors}
  \longrightarrow
  \text{lazy ROM state}
  \longrightarrow
  \text{interpretation functions on ROM outputs}.
\]
For example, \ec{message\_hash\_query}, \ec{challenge\_hash\_query},
\ec{sampler\_expand\_query}, and \ec{matrix\_expand\_query} prevent accidental
identification of distinct hash domains.  \ec{HAETAE\_RO.FRO.get} supplies the
lazy random-oracle semantics.  Operations such as \ec{ro\_message\_hash} and
\ec{ro\_signing\_coins} interpret the raw oracle output as a typed object when a
wrapper consumes that output.  In the current scheme wrapper, the sampler output
is consumed as signing coins, while the matrix, message-hash, and challenge ROM
calls are proof-visible instrumentation around deterministic internal
operations.

This is modeled hash-interface support, not a standard-model SHAKE theorem.
\ec{HAETAE\_FIPS202.ec} and \ec{HAETAE\_Keccak1600.ec} are useful for specifying
or validating concrete hash components, but the provable-security theorem proved
in the dissertation is a random-oracle-model theorem for the formal HAETAE
scheme.

\section{Residual model/spec obligations}
\label{sec:model-spec-residual-obligations}

The following obligations remain outside the central EasyCrypt security theorem
and should be tracked whenever the model or the specification changes.

\begin{longtable}{p{0.25\textwidth}p{0.65\textwidth}}
\caption{Model/spec obligations outside the central game proof.}
\label{tab:model-spec-obligations}\\
\toprule
\textbf{Obligation} & \textbf{Reason} \\
\midrule
\endfirsthead
\toprule
\textbf{Obligation} & \textbf{Reason} \\
\midrule
\endhead
Internal algorithm fidelity & \ec{keygen\_internal}, \ec{sign\_internal}, and
\ec{verify\_internal} must be audited against Figure 8 and the algebraic
specification.  The ROM wrapper exposes oracle structure but does not by itself
prove that every internal operation is the same as the specification formula. \\
\midrule
Deterministic sampling path & The specification derives signing randomness from
secret material and the message digest.  The proof model uses typed randomness
and ROM-mediated sampler expansion.  The correspondence is the explicit
modeling assumption in
\Cref{assm:signing-randomness-modeling-bridge}, not a checked theorem. \\
\midrule
Concrete hash instantiation & The checked implication is in the ROM.  Replacing
the ROM by SHAKE is a computational-assumption or indifferentiability-style step
not proved by the game theorem itself. \\
\midrule
Encoding and parsing & The proof-level signature type abstracts over concrete
byte encodings.  Implementation-level validity of decoding, rejection, and
malformed-input behavior is a separate verification target. \\
\midrule
Runtime overhead & The checked implication accounts for probabilities and query
budgets.  Concrete running time and circuit-size overhead of reductions are
recorded as an external ledger unless a future cost-aware formalization is added. \\
\bottomrule
\end{longtable}

This chapter does not weaken the machine-checked implication.  It states
precisely what the checked surface proves: a conditional ROM security result for
the formal HAETAE model, together with an explicit audit trail showing how that
model corresponds to the February 2026 specification.
